\begingroup
\setlength{\parindent}{0pt}
\setlength{\parskip}{3pt}
\setlength{\abovedisplayskip}{5pt plus 1pt minus 1pt}
\setlength{\belowdisplayskip}{5pt plus 1pt minus 1pt}
\setlength{\abovedisplayshortskip}{3pt plus 1pt}
\setlength{\belowdisplayshortskip}{5pt plus 1pt minus 1pt}
\clubpenalty=10000
\widowpenalty=10000
\displaywidowpenalty=10000

\subsection{Суммы последовательностей}
Сумма геометрической прогрессии:
\[
  S_n = \begin{cases}
    \dfrac{b_1(1 - q^n)}{1 - q}, & q \ne 1, \\
    nb_1, & q = 1.
  \end{cases}
\]
Сумма бесконечной геометрической прогрессии:
\[
  S = \frac{b_1}{1 - q}, \qquad |q| < 1.
\]
Сумма арифметической прогрессии:
\[
  \begin{aligned}
    S_n &= \frac{a_1 + a_n}{2} \cdot n \\
        &= \frac{2a_1 + d(n - 1)}{2} \cdot n.
  \end{aligned}
\]
\[
  \begin{aligned}
    1^2 + 2^2 + \dots + n^2 &= \frac{n(n + 1)(2n + 1)}{6}, \\
    1^3 + 2^3 + \dots + n^3 &= \frac{n^2(n + 1)^2}{4}.
  \end{aligned}
\]

\subsection{Вычисления}
Площадь криволинейного сектора:
\[
  \mu(F) = \frac{1}{2} \int_{\alpha}^{\beta} r^2(\varphi)\,d\varphi.
\]
Интегралы (при $a > 0$, на интервалах, где выражения определены):
\[
  \begin{aligned}
    \int \sqrt{a^2 - x^2}\,dx
      &= \frac{x}{2}\sqrt{a^2 - x^2} \\
      &\quad + \frac{a^2}{2}\arcsin\frac{x}{a} + C.
  \end{aligned}
\]
\[
  \begin{aligned}
    \int \sqrt{x^2 \pm a^2}\,dx
      &= \frac{x}{2}\sqrt{x^2 \pm a^2} \\
      &\quad \pm \frac{a^2}{2}\ln\left|x + \sqrt{x^2 \pm a^2}\right| + C.
  \end{aligned}
\]
\[
  \int \frac{dx}{a^2 + x^2}
    = \frac{1}{a}\operatorname{arctg}\frac{x}{a} + C.
\]
\[
  \int \frac{dx}{a^2 - x^2}
    = \frac{1}{2a}\ln\left|\frac{a + x}{a - x}\right| + C.
\]
\[
  \int \frac{x\,dx}{a^2 \pm x^2}
    = \pm\frac{1}{2}\ln|a^2 \pm x^2| + C.
\]
\[
  \int \frac{dx}{\sqrt{a^2 - x^2}}
    = \arcsin\frac{x}{a} + C.
\]
\[
  \int \frac{dx}{\sqrt{x^2 \pm a^2}}
    = \ln\left|x + \sqrt{x^2 \pm a^2}\right| + C.
\]
\[
  \int \frac{x\,dx}{\sqrt{a^2 \pm x^2}}
    = \pm\sqrt{a^2 \pm x^2} + C.
\]

% Keep combinatorics together in the second column.
\newpage
\subsection{Комбинаторика}
Размещения:
\begin{itemize}
  \item Без повторений: $A_n^k = \dfrac{n!}{(n - k)!}$.
  \item С повторениями: $\overline{A}_n^k = n^k$.
\end{itemize}
Сочетания:
\begin{itemize}
  \item Без повторений: $C_n^k = \dfrac{n!}{k!(n - k)!}$.
  \item С повторениями: $\overline{C}_n^k = \dfrac{(n + k - 1)!}{k!(n - 1)!}$.
\end{itemize}
Перестановки:
\begin{itemize}
  \item Без повторений: $P_n = n!$.
  \item С повторениями:
    \[
      P_n(n_1, n_2, \dots, n_k) = \frac{n!}{n_1!n_2!\dots n_k!}.
    \]
\end{itemize}
\[
  \begin{aligned}
    C_n^k &= C_{n - 1}^k + C_{n - 1}^{k - 1}, \\
    C_n^k &= C_n^{n - k}.
  \end{aligned}
\]

\subsubsection*{Stars and bars}
При $n \ge 0$, $k \ge 1$ число представлений $n$ в виде суммы $k$ неотрицательных слагаемых:
\[
  C_{n + k - 1}^{k - 1}.
\]
При $n \ge k \ge 1$ число представлений $n$ в виде суммы $k$ положительных слагаемых:
\[
  C_{n - 1}^{k - 1}.
\]

% Reference checks: cp-algorithms.com/combinatorics/binomial-coefficients.html,
% cp-algorithms.com/combinatorics/burnside.html, dlmf.nist.gov/26.8.
\subsubsection*{Биномиальные тождества}
При $n \ge 0$ считаем $C_n^k=0$ для $k<0$ и $k>n$.
Для неотрицательных целых параметров:
\[
  \begin{gathered}
    \sum_{k=0}^{n} C_n^k=2^n, \qquad
    \sum_{k=0}^{n}(C_n^k)^2=C_{2n}^n, \\
    \sum_{k=0}^{n} kC_n^k=n2^{n-1}\quad(n\ge1), \\
    \sum_{k=0}^{m} C_{r+k}^{r}=C_{r+m+1}^{r+1}, \qquad
    \sum_{k=0}^{n} C_r^k C_s^{n-k}=C_{r+s}^{n}.
  \end{gathered}
\]

\textbf{Числа Каталана:}
\[
  \mathrm{Cat}_n=\frac{C_{2n}^{n}}{n+1}
    =C_{2n}^{n}-C_{2n}^{n+1},\qquad \mathrm{Cat}_0=1.
\]
Считают правильные скобочные последовательности из $n$ пар и триангуляции выпуклого $(n+2)$-угольника ($n\ge1$).

\textbf{Беспорядки} (перестановки без неподвижных точек):
\[
  \begin{gathered}
    D_0=1,\qquad D_1=0,\\
    D_n=(n-1)(D_{n-1}+D_{n-2})\quad(n\ge2).
  \end{gathered}
\]

\Needspace{12\baselineskip}
\textbf{Числа Стирлинга.} $S(n,k)$ считает разбиения $n$ различных элементов на $k$ непустых неупорядоченных групп; $c(n,k)$ считает перестановки с $k$ циклами:
\[
  \begin{aligned}
    S(n,k)&=S(n-1,k-1)+kS(n-1,k),\\
    c(n,k)&=c(n-1,k-1)+(n-1)c(n-1,k).
  \end{aligned}
\]
Обе базы: значение в $(0,0)$ равно $1$, в $(n,0)$ и $(0,n)$ при $n>0$ равно $0$. Число сюръекций на $k$ различных значений равно $k!S(n,k)$.

\subsubsection*{Включения-исключения и симметрии}
Для $m$ множеств, где $[m]=\{1,\dots,m\}$:
\[
  \left|\bigcup_{i=1}^{m} A_i\right|
  =\sum_{\varnothing\ne S\subseteq[m]}
    (-1)^{|S|+1}\left|\bigcap_{i\in S}A_i\right|.
\]
Число решений $x_1+\dots+x_k=n$, $0\le x_i\le b$, для $n,b\ge0$, $k\ge1$:
\[
  \sum_{j=0}^{\min(k,\lfloor n/(b+1)\rfloor)}
    (-1)^j C_k^j C_{n-j(b+1)+k-1}^{k-1}.
\]
\textbf{Лемма Бёрнсайда.} Для конечной группы симметрий $G$ число классов эквивалентности:
\[
  \frac{1}{|G|}\sum_{g\in G}|\operatorname{Fix}(g)|.
\]
$\operatorname{Fix}(g)$ --- объекты, не меняющиеся при $g$. При независимой раскраске позиций в $c$ цветов перестановка с $q$ циклами сохраняет $c^q$ раскрасок. Ожерелья из $n\ge1$ бусин с точностью до поворота (без отражений):
\[
  \frac1n\sum_{i=0}^{n-1} c^{\gcd(n,i)}
  =\frac1n\sum_{d\mid n}\varphi(d)c^{n/d}.
\]
Деление в целочисленной формуле нельзя автоматически заменять умножением на обратное по модулю: обратное должно существовать.

% Reference checks: cp-algorithms.com/algebra/chinese-remainder-theorem.html,
% dlmf.nist.gov/27.5; the CRT merge below follows directly from x = r1 + m1*t.
\subsection{Теория чисел и модульная арифметика}
\textbf{Диофантово уравнение.} Пусть $g=\gcd(a,b)>0$, $ax_0+by_0=g$. Уравнение $ax+by=c$ разрешимо тогда и только тогда, когда $g\mid c$. Все решения, $t\in\mathbb Z$:
\[
  x=x_0\frac cg+\frac bg t,\qquad
  y=y_0\frac cg-\frac ag t.
\]
Для таких $a,b$: $\operatorname{lcm}(a,b)=|(a/g)b|$.

\textbf{Обратный элемент.} При $m>1$ число $a^{-1}\bmod m$ существует ровно при $\gcd(a,m)=1$: это коэффициент при $a$ в \texttt{egcd}, нормализованный в $[0,m)$.
При взаимной простоте $a^{\varphi(m)}\equiv1\pmod m$.
Для простого $p$ и $p\nmid a$: $a^{-1}\equiv a^{p-2}\pmod p$.
Сокращать показатель по $\varphi(m)$ можно при $\gcd(a,m)=1$.
Сравнение $ax\equiv b\pmod m$ имеет решения ровно при $\gcd(a,m)\mid b$; тогда их $\gcd(a,m)$ по модулю $m$.

\textbf{Китайская теорема об остатках.} Объединяем $x\equiv r_1\pmod{m_1}$ и $x\equiv r_2\pmod{m_2}$, $m_1,m_2>0$. Пусть $g=\gcd(m_1,m_2)$, $\Delta=r_2-r_1$, $q=m_2/g$. Решение существует ровно при $g\mid\Delta$:
\[
  \begin{aligned}
    t&\equiv \frac{\Delta}{g}\left(\frac{m_1}{g}\right)^{-1}\pmod q,\\
    x&\equiv r_1+m_1t\pmod{m_1q}.
  \end{aligned}
\]
При $q=1$ берём $t=0$. Ответ единственен по модулю $\operatorname{lcm}(m_1,m_2)=m_1q$. Промежуточные произведения при необходимости считать в \texttt{\_\_int128}.

\textbf{Мультипликативные функции.} Для $n=\prod p_i^{a_i}\ge1$:
\[
  \begin{aligned}
    \tau(n)&=\prod_i(a_i+1),&
    \sigma(n)&=\prod_i\frac{p_i^{a_i+1}-1}{p_i-1},\\
    \varphi(n)&=n\prod_{p\mid n}\left(1-\frac1p\right),&
    \sum_{d\mid n}\varphi(d)&=n.
  \end{aligned}
\]
$\tau$ --- число делителей, $\sigma$ --- их сумма.
Функция Мёбиуса: $\mu(1)=1$; $\mu(n)=0$, если $p^2\mid n$ для некоторого простого $p$; иначе $\mu(n)=(-1)^r$, где $r$ --- число различных простых делителей.
\[
  \begin{gathered}
    \sum_{d\mid n}\mu(d)=[n=1],\\
    F(n)=\sum_{d\mid n}f(d)
    \ \Longleftrightarrow\ f(n)=\sum_{d\mid n}\mu(d)F(n/d).
  \end{gathered}
\]
Здесь $[P]$ равно $1$, если $P$ истинно, иначе $0$.
Число пар $1\le a\le A$, $1\le b\le B$ с $\gcd(a,b)=1$:
\[
  \sum_{d=1}^{\min(A,B)}\mu(d)
    \left\lfloor\frac Ad\right\rfloor
    \left\lfloor\frac Bd\right\rfloor.
\]

\textbf{Факториалы по простому модулю.} Степень простого $p$ в $n!$:
\[
  v_p(n!)=\sum_{j\ge1}\left\lfloor\frac{n}{p^j}\right\rfloor.
\]
\textbf{Люка:} если $n_i,k_i$ --- цифры $n,k$ в системе счисления по простому основанию $p$, то
$C_n^k\equiv\prod_i C_{n_i}^{k_i}\pmod p$.
Формула через факториалы и обратные факториалы напрямую работает при $0\le k\le n<p$.

\textbf{Группировка частных.} Для $n\ge1$, $1\le l\le n$ значение $q=\lfloor n/l\rfloor$ постоянно на $[l,\lfloor n/q\rfloor]$. Переход к $l=\lfloor n/q\rfloor+1$ перебирает все такие блоки за $O(\sqrt n)$.

\textbf{Монеты.} Для взаимно простых $a,b\ge2$ наибольшее целое, не представимое как $ax+by$, $x,y\ge0$, равно $ab-a-b$; число положительных непредставимых равно $(a-1)(b-1)/2$.

% Reference checks: cses.fi/book/book.pdf, chapter 20 (Hall, Konig, DAG cover);
% cp-algorithms.com/graph/kirchhoff-theorem.html.
\subsection{Графы и деревья}
\textbf{Эйлеров путь} проходит по каждому ребру ровно один раз. При наличии рёбер вершины ненулевой степени должны быть связны после забывания направлений.
В неориентированном графе дополнительно нужно $0$ или $2$ вершины нечётной степени; цикл существует при $0$.
В ориентированном графе для цикла все $\deg^+=\deg^-$; для незамкнутого пути у начала $\deg^+-\deg^-=1$, у конца $-1$, у остальных $0$.

\textbf{Паросочетания.} В двудольном графе размер максимального паросочетания равен размеру минимального вершинного покрытия (Кёниг). Покрытие касается каждого ребра. Дополнение минимального покрытия --- максимальное независимое множество.
По Холлу паросочетание покрывает всю левую долю $L$ ровно тогда, когда $|N(S)|\ge|S|$ для каждого $S\subseteq L$.

\textbf{Покрытие DAG путями.} Минимальное число непересекающихся по вершинам путей, покрывающих все $n$ вершин, равно $n-\nu$. Здесь $\nu$ --- максимум паросочетания между двумя копиями вершин с рёбрами $u_Lv_R$ для каждого исходного $u\to v$. Одновершинные пути разрешены.

\textbf{Потоки и разрезы.} Максимальный поток равен минимальной пропускной способности $s$--$t$-разреза. При единичных ёмкостях он равен максимуму непересекающихся по рёбрам $s$--$t$-путей. Для ограничения использования вершины разделяем её на вход и выход, соединённые ребром ёмкости $1$; это ограничение не накладываем на $s,t$.

\textbf{MST.} В связном неориентированном взвешенном графе единственное минимальное ребро разреза входит во все минимальные остовные деревья; единственное максимальное ребро цикла не входит ни в одно.

\textbf{2-SAT.} Клауза $(a\lor b)$ даёт импликации $\neg a\to b$, $\neg b\to a$. Формула выполнима ровно тогда, когда ни для одной переменной $x$ вершины $x,\neg x$ не лежат в одной компоненте сильной связности графа импликаций.

\textbf{Деревья.} Для невзвешенного дерева диаметр находится двумя поисками самой далёкой вершины. Если $A,B$ --- его концы, то
$\max_u d(v,u)=\max(d(v,A),d(v,B))$.
Радиус равен $\lceil d(A,B)/2\rceil$; центры --- одна или две средние вершины диаметра.
\emph{Центроид} при удалении оставляет компоненты размера не более $n/2$; центроидов один или два, во втором случае они соседние.

\textbf{Кэли и Прюфер.} Помеченных деревьев на $n\ge2$ вершинах ровно $n^{n-2}$. Им взаимно однозначно соответствуют последовательности Прюфера длины $n-2$; $\deg(v)=1+$ число вхождений $v$ в последовательность.

\textbf{Кирхгоф.} Для неориентированного графа без петель пусть $L_{ii}=\deg(i)$, $L_{ij}=-m_{ij}$ при $i\ne j$, где $m_{ij}$ --- число рёбер между $i,j$. Удаляем одну и ту же строку и столбец из $L$: определитель оставшейся матрицы равен числу остовных деревьев.

\textbf{Планарность.} В плоской укладке графа с $c$ компонентами $n-m+f=1+c$, включая внешнюю грань. Для простого планарного графа при $n\ge3$: $m\le3n-6$, а для двудольного $m\le2n-4$.

\subsection{Вероятности и математическое ожидание}
\textbf{Линейность:} $\mathbb E[\sum_i X_i]=\sum_i\mathbb E[X_i]$; независимость не нужна. Для индикатора события $A$: $\mathbb E[\mathbf1_A]=P(A)$.
При конечных вторых моментах:
\[
  \operatorname{Var}(X)=\mathbb E[X^2]-(\mathbb E[X])^2.
\]
Дисперсии независимых величин складываются.
Для целочисленного $X\ge0$: $\mathbb E[X]=\sum_{k\ge1}P(X\ge k)$.

\textbf{Первый успех.} При независимых попытках с вероятностью успеха $p>0$ число попыток до первого успеха включительно имеет ожидание $1/p$ и дисперсию $(1-p)/p^2$.
Чтобы собрать все $n$ равновероятных типов при независимых выборах с возвращением, в среднем нужно $nH_n$ выборов, где $H_n=\sum_{k=1}^n1/k$.

\textbf{Первый шаг.} Ожидаемая суммарная стоимость до завершения:
\[
  E_v=c_v+\sum_u p_{vu}E_u,\qquad E_{\rm terminal}=0.
\]
$c_v$ --- стоимость шага, $p_{vu}$ --- вероятности переходов. В DAG считаем в обратном топологическом порядке; при циклах решаем линейную систему, предварительно проверив конечность ожидания.

\textbf{Оценки.} Для $t>0$:
\[
  \begin{aligned}
    P(X\ge t)&\le\mathbb E[X]/t &&(X\ge0),\\
    P(|X-\mathbb E[X]|\ge t)&\le\operatorname{Var}(X)/t^2.
  \end{aligned}
\]

% Reference check: cp-algorithms.com/game_theory/sprague-grundy-nim.html.
\subsection{Игры}
\textbf{Шпраг--Гранди.} Конечная игра без случайности и скрытой информации, допустимые ходы из позиции одинаковы для обоих игроков, граф ходов ацикличен; не имеющий хода проигрывает:
\[
  g(v)=\operatorname{mex}\{g(u):v\to u\},\qquad
  \operatorname{mex}\varnothing=0.
\]
$\operatorname{mex}$ --- наименьшее отсутствующее неотрицательное целое. В сумме независимых игр, где за ход меняется ровно одна компонента, число Гранди равно XOR чисел компонент; позиция проигрышная ровно при XOR $=0$.

\textbf{Ним.} За ход уменьшают одну кучу; $g(a)=a$. При ненулевом XOR $s$ нужно выбрать кучу $a$ с $a\mathbin{\oplus}s<a$ и заменить её размер на $a\mathbin{\oplus}s$.
В \emph{мизерном ниме} взявший последний камень проигрывает: если все непустые кучи единичные, выигрышен чётный их счёт; иначе действует обычный критерий XOR.
Если из одной кучи можно брать от $1$ до $k$ камней и последний ход выигрывает, проигрышны ровно размеры, кратные $k+1$.

% Reference check: cp-algorithms.com/geometry/picks-theorem.html.
\subsection{Геометрические формулы}
Для невырожденного треугольника со сторонами $a,b,c$, полупериметром $s=(a+b+c)/2$, площадью $S$, радиусами вписанной и описанной окружностей $r,R$:
\[
  \begin{gathered}
    S=\sqrt{s(s-a)(s-b)(s-c)}=rs=\frac{abc}{4R},\\
    c^2=a^2+b^2-2ab\cos C,\\
    \frac a{\sin A}=\frac b{\sin B}=\frac c{\sin C}=2R.
  \end{gathered}
\]
Углы $A,B,C$ лежат напротив соответствующих сторон.
Для прямой $ax+by+c=0$, $a^2+b^2>0$, расстояние от $(x_0,y_0)$ равно $|ax_0+by_0+c|/\sqrt{a^2+b^2}$.
Проекция точки $P$ на прямую $AB$, $v=B-A\ne0$:
$A+v\,((P-A)\cdot v)/(v\cdot v)$; для отрезка ограничить коэффициент при $v$ диапазоном $[0,1]$.

\textbf{Пик.} У простого многоугольника без дыр с целочисленными вершинами $S=I+B/2-1$, где $I$ --- число внутренних решётчатых точек, $B$ --- граничных.
Для циклически перечисленных вершин
$B=\sum_i\gcd(|x_{i+1}-x_i|,|y_{i+1}-y_i|)$.
На решётчатом отрезке с различными концами точек $\gcd(|\Delta x|,|\Delta y|)+1$.

\textbf{Манхэттенское расстояние.} При $u=x+y$, $v=x-y$:
\[
  |x_1-x_2|+|y_1-y_2|
  =\max(|u_1-u_2|,|v_1-v_2|).
\]
Для обратного перехода к целочисленным $x,y$ числа $u,v$ должны иметь одинаковую чётность.

% Reference checks: cp-algorithms.com/string/prefix-function.html,
% cp-algorithms.com/dynamic_programming/divide-and-conquer-dp.html,
% cp-algorithms.com/dynamic_programming/knuth-optimization.html.
\Needspace{16\baselineskip}
\subsection{Строки, алгебра и оптимизации}
\textbf{Период строки.} При $n>0$ и префикс-функции $\pi$ кратчайший период равен $p=n-\pi[n-1]$. Строка состоит из целого числа повторений блока длины $p$, если $n\bmod p=0$; иначе кратчайший блок для такого представления --- вся строка.
Длины всех непустых собственных границ (совпадающих префиксов и суффиксов): $k=\pi[n-1]$, затем $k=\pi[k-1]$ до нуля.

\textbf{Различные подстроки.} Если $\mathrm{LCP}_i$ --- длины общих префиксов соседних суффиксов в лексикографическом порядке, то их число равно
\[
  \frac{n(n+1)}2-\sum_{i=1}^{n-1}\mathrm{LCP}_i.
\]
Различных непустых подстрок-палиндромов не больше $n$.

\textbf{Линейная алгебра.} Над полем система $Ax=b$ совместна ровно при $\operatorname{rank}(A)=\operatorname{rank}(A\mid b)$. При $m$ неизвестных и ранге $r$ над $\mathbb F_p$ ($p$ простое) число решений совместной системы равно $p^{m-r}$.
Для $n$ двоичных векторов ранга $r$ существует $2^r$ различных XOR подмножеств; каждый достижимый XOR получается ровно $2^{n-r}$ способами, включая пустое подмножество.

\textbf{Матричные степени.} $(A^k)_{uv}$ матрицы смежности считает маршруты длины $k$ из $u$ в $v$ (вершины и рёбра могут повторяться). Линейная рекуррента порядка $d$ вычисляется матричным возведением в степень за $O(d^3\log n)$.

\textbf{Divide and conquer DP.} Для
\[
  dp_t[j]=\min_{0\le k<j}\bigl(dp_{t-1}[k]+C(k,j)\bigr)
\]
монотонность $opt_t[j]\le opt_t[j+1]$ даёт $O(n\log n)$ на слой при $O(1)$ на стоимость. Достаточно неравенства Монжа
$C(a,c)+C(b,d)\le C(a,d)+C(b,c)$ при $a\le b<c\le d$.

\textbf{Кнут (отрезочная DP).} При $dp[i][i]=0$ и
\[
  dp[i][j]=C(i,j)+\min_{i\le k<j}
    \bigl(dp[i][k]+dp[k+1][j]\bigr)
\]
условие $opt[i][j-1]\le opt[i][j]\le opt[i+1][j]$ даёт $O(n^2)$ при $O(1)$ на стоимость.
Достаточно $C(b,c)\le C(a,d)$ и
$C(a,c)+C(b,d)\le C(a,d)+C(b,c)$ для $a\le b\le c\le d$.
При равных минимумах в обеих оптимизациях последовательно выбираем наибольший индекс.

\textbf{Полезные оценки и экстремумы.}
Все пары «маска, её подмаска» для $n$ бит: $3^n$.
Суммы по подмаскам для всех масок (SOS DP): $O(n2^n)$.
Минимальное число невозрастающих подпоследовательностей, на которые разбивается последовательность, равно длине её строго возрастающей НВП.
Сумма $\sum_i|a_i-x|$ минимальна на медиане, а $\sum_i(a_i-x)^2$ --- на среднем арифметическом.

% Reference checks: cp-algorithms.com/algebra/fibonacci-numbers.html,
% dlmf.nist.gov/3.3; atcoder.github.io/ac-library/production/document_en/convolution.html
% and atcoder/internal_math.hpp (primitive_root_constexpr).
\subsection{Дополнительные алгебраические факты}
\textbf{Фибоначчи.} $F_0=0$, $F_1=1$, $F_{n+2}=F_{n+1}+F_n$.
Быстрое удвоение вычисляет пару $(F_n,F_{n+1})$ за $O(\log n)$ арифметических операций:
\[
  \begin{aligned}
    F_{2k}&=F_k(2F_{k+1}-F_k),\\
    F_{2k+1}&=F_k^2+F_{k+1}^2.
  \end{aligned}
\]
Также $\gcd(F_a,F_b)=F_{\gcd(a,b)}$ и $\sum_{i=0}^n F_i=F_{n+2}-1$ для неотрицательных индексов.

\textbf{Перестановки.} Если в перестановке $n$ элементов $c$ циклов, включая неподвижные точки, то её знак $(-1)^{n-c}$, а минимальное число обменов произвольных пар для сортировки равно $n-c$.
Для соседних обменов минимум равен числу инверсий.
Порядок перестановки (число применений до тождественной) равен НОК длин циклов.

\textbf{Интерполяция Лагранжа.} Над полем $n+1$ точка $(x_i,y_i)$ с попарно различными $x_i$ задаёт единственный многочлен степени не выше $n$:
\[
  P(x)=\sum_{i=0}^{n} y_i
    \prod_{\substack{0\le j\le n\\j\ne i}}
      \frac{x-x_j}{x_i-x_j}.
\]
По простому модулю $p$ узлы должны быть различны именно по модулю $p$. Для узлов $0,\dots,n$, $n<p$, значение в одной точке считается за $O(n)$ после подготовки факториалов и обратных факториалов: числители --- префиксные и суффиксные произведения, знаменатель $(-1)^{n-i}i!(n-i)!$.

\textbf{NTT.} Для точного преобразования по простому модулю $p$ длины $N=2^k$ нужно $N\mid(p-1)$. Если $g$ --- первообразный корень, то $w=g^{(p-1)/N}$ имеет порядок $N$.
Обратное преобразование использует $w^{-1}$ и множитель $N^{-1}$.
Для свёртки двух массивов выбираем $N\ge |a|+|b|-1$.
Частый модуль: $998244353=119\cdot2^{23}+1$, первообразный корень $g=3$, допустимы $N\le2^{23}$.
\endgroup
